\section{Introduction}

% \paragraph{Skills as a practical interface for agent behavior.}
Agentic large language models increasingly rely on external skills~\citep{xu2026agent,jiang2026sok,chen2026cua,nguyen2025guisurvey,huang2025towardsagentic}, 
which are reusable natural-language procedures that specify how an agent should decompose a task, invoke tools, and verify intermediate results. 
Since skills are designed as complementary modules of the task language model, 
they provide a practical interface for shaping LLM behavior without updating the task language model. 
This separation is practical in realistic deployments, where the task model may be proprietary~\citep{gao2025survey,huang2025surveyfoundation}, 
expensive to adapt~\citep{zhou2025memento,wu2024personalizedmllm,wu2025mitigatingforget,wang2025selfupdatable,zhang2024personalizationsurvey}, or shared across many downstream tasks~\citep{li2026organizing,chen2026skillcraft,jiang2026xskill}.

% \paragraph{The challenge of recurrent instance-level skill improvement.}
However, many skill-based systems rely on well-engineered skills as predefined and fixed artifacts, 
or revise them only through prompting and single-round self-refinement~\citep{jiang2026sok,xu2026agent,chen2026cua,jiang2026xskill,chen2026skillcraft,xia2026skillrl,zhang2026memskill}. 
These strategies can be misaligned and inefficient, since they rely on repeated prompt edits and are not directly aligned with downstream rewards~\citep{wu2025incontext,kveton2025activedpo,xia2025surveyactive}.
On the other hand, some recent works adapt the task language model to the skill by fine-tuning the task language model on the skill~\citep{wang2025reinforcement,li2025commit,wang2024instructgraph},
but this approach can be computationally costly and limited in adaptivity to many downstream tasks~\citep{zhou2025memento,xia2026skillrl,jiang2026adaptationagenticaisurvey,wu2025mitigatingforget,yao2024federatedllm}.
Moreover, skill optimization should not be a single prompt-editing step, but a recurrent decision problem: 
a useful skill should improve rollout quality under the current context, 
and a useful revision should transform observed successes and failures into a better skill for the next round. 
This naturally leads to a bi-level optimization objective that couples execution quality within each generation with skill improvement across generations.
% \tong{Might it be helpful to explain in the introduction why existing single-level optimization approaches (e.g., standard RL or GRPO) are insufficient for this recurrent skill optimization problem?}

% \paragraph{Overview of Skill-R1.}
To address this problem, we propose \textbf{Skill-R1}, a reinforcement learning framework for recurrent agent skill evolution from verifiable rewards. 
\textbf{Skill-R1} separates \emph{task execution} from \emph{skill improvement}. 
A frozen task LLM produces rollouts conditioned on the current skill, 
while a lightweight skill generator is trained to revise that skill based on the task context, prior rollouts, and their verified outcomes~\citep{xia2025knowledgequery,wang2025dice,vannguyen2025smallsurvey}. 
By updating only the skill generator, the framework remains compatible with both open- and closed-source task models and can be more efficient than updating the task model itself~\citep{vannguyen2025smallsurvey,xia2025surveyactive}.

% \paragraph{Evolutionary rollouts across generations.}
The core interaction loop of \textbf{Skill-R1} is multi-generation. 
For a given instance, the current skill induces a group of rollouts from the frozen task model, a verifier scores these rollouts, and the resulting trials and errors~\citep{wu2025ocean, wu2024decot},
and reward signals are appended to an instance-specific history that conditions the next skill revision. 
Recurring execution of this process yields an evolutionary view of skill improvement, in which each generation proposes a new skill, 
tests it through a rollout population, and passes forward evidence for subsequent refinement.
This recurrent structure makes it possible to optimize directional improvement~\citep{xia2026skillrl,sun2025seagent,zhou2025memento} on the same instance rather than selecting among isolated one-shot attempts.

% \paragraph{Bi-level credit assignment for skill evolution.}
The policy optimization of the skill generator requires both intra-generation and inter-generation credit assignment to optimize the skill generator.
To enable optimization of the recurrent skill generator, we introduce a bi-level group-relative policy optimization~\citep{zhong2026rc,shao2024deepseekmath,mundada2026wsgrpo} objective that mirrors the recurrent structure above. 
The \emph{intra-generation} term compares rollouts sampled under the same skill and captures relative execution quality within a generation. 
The \emph{inter-generation} term measures whether a revised skill improves the average verified performance relative to the previous generation.
Together, these signals provide a principled objective for rewarding both strong rollouts and beneficial skill revisions~\citep{wu2025incontext, surana2026massdpo, li2025impsamplingdpo, huang2025pluralistic}.
% \fixme{Empirical findings}
Our experiments demonstrate that Skill-R1 consistently improves agent performance across diverse reasoning and tool-use benchmarks, outperforming both no-skill baselines and standard GRPO. 
We also show a clear performance gap between skill generations, which demonstrates the effectiveness of the recurrent skill evolution.
% Instead of relying on one-shot edits, Skill-R1 adopts an iterative evolution process guided by execution feedback, enabling more reliable and targeted improvements. The multi-generation framework supports progressive capability building, while the bi-level objective ensures that updates are both locally effective and aligned with longer-term optimization. 
% Furthermore, learning the skill editor plays a critical role by allowing the system to distinguish persistent failure modes from incidental noise, resulting in more stable and transferable skill adaptations.
We summarize the contributions of this work as follows:
\begin{itemize}
    \item We formulate recurrent instance-level skill evolution as a bi-level group-relative policy optimization problem.
    \item We propose \textbf{Skill-R1}, a model-agnostic framework that trains a lightweight skill generator while keeping the task LLM frozen.
    \item We introduce a recurrent multi-generation rollout process together with a bi-level GRPO objective that combines \emph{intra-} and \emph{inter-generation} advantages. 
    % \tong{Might it be helpful to verify that the benefit of combining intra- and inter-generation advantages is demonstrated through ablation studies in the experimental section?}
    \item We evaluate \textbf{Skill-R1} on agent tasks with verifiable rewards, comparing with agent skill and GRPO baselines and achieving superior performance.
    % comparing against fixed-skill and single-generation baselines and isolating the effects of multi-generation evolution and bi-level advantages through ablations.} \tong{It seems that the keyword 'ablation' is only mentioned here in this paper.}
\end{itemize}